\documentclass{article}
\usepackage{graphicx}
\usepackage[margin=0.75in]{geometry}

% code styling
\usepackage{xcolor}
\usepackage{listings}

\lstdefinelanguage{Golang}{
  morekeywords={
    break,default,func,interface,select,case,defer,go,map,struct,
    chan,else,goto,package,switch,const,fallthrough,if,range,type,
    continue,for,import,return,var,bool,string,int,int8,int16,int32,
    int64,uint,uint8,uint16,uint32,uint64,uintptr,byte,rune,
    float32,float64,complex64,complex128,nil,true,false,iota
  },
  sensitive=true,
  morecomment=[l]{//},
  morecomment=[s]{/*}{*/},
  morestring=[b]",
  morestring=[b]',
  morestring=[s]{`}{`}
}

\lstdefinelanguage{Rust}{
  morekeywords={
    abstract,alignof,as,become,break,const,continue,crate,do,
    else,enum,extern,false,fn,for,if,impl,in,let,loop,match,mod,move,
    mut,pub,ref,return,self,Self,static,struct,super,trait,true,type,
    typeof,unsafe,unsized,use,where,while,yield,str,String,Vec,Option,Result
  },
  otherkeywords={&, !, ', ->, =>},
  sensitive=true,
  morecomment=[l]{//},
  morecomment=[s]{/*}{*/},
  morestring=[b]",
  moredelim=[s][\color{orange!70!black}]{\ '}{\ } % Highlight lifetimes like 'a
}

\lstset{
    language=C++,
    basicstyle=\ttfamily\large,
    keywordstyle=\color{blue}\bfseries,
    stringstyle=\color{red!70!black},
    commentstyle=\color{green!50!black}\itshape,
    numberstyle=\normalsize\color{gray},
    numbers=left,
    stepnumber=1,
    numbersep=8pt,
    frame=single,
    rulecolor=\color{black!30},
    breaklines=true,
    tabsize=4,
    showstringspaces=false,
    backgroundcolor=\color{gray!5},
    % C++ Specific Keywords
    emph={std, ptrdiff_t, size_t, uint64_t, int64_t, atomic, mutex, condition_variable, barrier, counting_semaphore},
    emphstyle=\color{teal}\bfseries
}

\title{CS3211 Technical Reference}
\author{ngmh}
\date{04/05/2026}

\begin{document}
\large
\maketitle

\section{C++ Threads}
\begin{lstlisting}[language=C++]
// function
void hello() {
    std::cout << "Hello";
}

// function object
class task {
public:
    void operator()() const {
        hello();
    }
}

int main() {
    std::thread t(hello);
    t.join();

    task f;
    std::thread my_thread(f);
}

\end{lstlisting}

\subsection{Most Vexing Parse}
The following code defines a function that:
\begin{itemize}
    \item Takes in a pointer to a function with no parameters and returns a task
    \item Returns a thread object
\end{itemize}
\begin{lstlisting}[language=C++]
std::thread my_thread(task());
\end{lstlisting}

Instead, to create a thread, it should be:
\begin{lstlisting}[language=C++]
std::thread my_thread((task()));
std::thread my_thread{task()};
std::thread my_thread([]{ hello(); });
\end{lstlisting}

\subsection{Waiting and Detaching}
Always join even when an exception occurs! If you forget to join a joinable thread, standard says \verb|std::terminate| is called.
\begin{lstlisting}[language=C++]
t.join();
t.joinable(); // true

t.detach();
t.joinable(); // false
\end{lstlisting}

\subsection{Passing Arguments to Threads}
Pass in references if necessary.
\begin{lstlisting}[language=C++]
void f(int i);
std::thread t(f, 3);

void update(widget w, data& data);
std::thread t(update, w, std::ref(data));
// necessary! constructor copies / moves objects
\end{lstlisting}

\subsection{Thread Ownership}
\begin{lstlisting}[language=C++]
// out of function
std::thread f() {
    void foo();
    return std::thread{foo};
}

// into function
void f(std::thread t)
void g() {
    void func();
    f(std::thread{func});
    // or
    std::thread t{func};
    f(std::move(t));
}

\end{lstlisting}

\section{C++ Synchronisation}

\subsection{RAII}
\begin{itemize}
    \item Resource Acquisition Is Initialisation
    \item Binds life cycle of resource that must be acquired before use to the lifetime of an object
    \item Resource is released when lifetime of controlling object ends in reverse order of acquisition
\end{itemize}

\subsection{Using Mutexes}
\begin{lstlisting}[language=C++]
std::mutex mut;
void f() {
    std::lock_guard<std::mutex> guard{mut};
    // critical section
}
\end{lstlisting}

\subsection{Types of Locks}
\begin{lstlisting}[language=C++]
std::mutex m1, m2;

void f() {
    // lock_guard is the simplest
    // can transfer ownership of lock
    std::lock(m1, m2);
    std::lock_guard<std::mutex> lock1(m1, std::adopt_lock);
    std::lock_guard<std::mutex> lock2(m2, std::adopt_lock);

    // unique lock ownership can be transferred
    // can manually lock and unlock with unique lock
    std::unique_lock<std::mutex> u_lock{m1, std::defer_lock};
    u_lock.lock();
    u_lock.unlock();
    // or
    std::lock(u_lock);
    
    // scoped lock can take in multiple mutexes, ensures no deadlock
    std::scoped_lock<std::mutex, std::mutex> guard{m1, m2};
}
\end{lstlisting}

\subsection{Condition Variables}
\begin{lstlisting}[language=C++]
std::mutex mut;
std::queue<int> q;
std::condition_variable data_cond;

void producer() {
    {
        std::lock_guard<std::mutex> lk{mut};
        q.push(1);
    }
    data_cond.notify_one();
}

void consumer() {
    while (true) {
        // condition_variable needs unique_lock
        // condition_variable_any does not
        std::unique_lock<std::mutex> lk{mut};
        data_cond.wait(lk, []{ return !q.empty(); });
        // data_cond.wait_for(...)
        // data_cond.wait_until(...)
        q.pop();
    }
}
\end{lstlisting}

\section{C++ Memory Orders}

\subsection{\texttt{std::memory\_order} Enumeration}
\begin{lstlisting}[language=C++]
// Default
std::memory_order_seq_cst
// Store Operations
std::memory_order_seq_cst
std::memory_order_release
std::memory_order_relaxed
// Load Operations
std::memory_order_seq_cst
std::memory_order_acquire
std::memory_order_consume
std::memory_order_relaxed
// RMW Operations
std::memory_order_seq_cst
std::memory_order_acq_rel
std::memory_order_acquire
std::memory_order_release
std::memory_order_consume
std::memory_order_relaxed
\end{lstlisting}

\subsection{Memory Orderings}
\begin{lstlisting}[language=C++]
// Sequentially Consistent Ordering
std::memory_order_seq_cst

// Relaxed Ordering
std::memory_order_relaxed

// Acquire-Release Ordering
std::memory_order_acquire
std::memory_order_release
std::memory_order_acq_rel
std::memory_order_consume
\end{lstlisting}

\section{Concurrent Data Structures}

\subsection{S1: Thread Safe Stack with one global mutex}
\begin{itemize}
    \item Safe for multiple threads to call member functions concurrently
    \item Exception safe
    \item Work is serialised; only one thread can ever do work, limiting performance
    \item No means of waiting for item to be added
\end{itemize}
\begin{lstlisting}[language=C++]
struct empty_stack: std::exception {
    const char* what() const throw();
};
template<typename T>
class threadsafe_stack {
private:
    std::stack<T> data;
    mutable std::mutex m;
public:
    threadsafe_stack(){}
    threadsafe_stack(const threadsafe_stack& other) {
        std::lock_guard<std::mutex> lock(other.m);
        data = other.data;
    }
    threadsafe_stack& operator=(const threadsafe_stack&) = delete;
    void push(T new_value) {
        std::lock_guard<std::mutex> lock(m);
        data.push(std::move(new_value));
    }
    std::shared_ptr<T> pop() {
        std::lock_guard<std::mutex> lock(m);
        if (data.empty()) throw empty_stack();
        std::shared_ptr<T> const res(
            std::make_shared<T>(std::move(data.top()))
        );
        data.pop();
        return res;
    }
    void pop(T& value) {
        std::lock_guard<std::mutex> lock(m);
        if (data.empty()) throw empty_stack();
        value = std::move(data.top());
        data.pop();
    }
    bool empty() const {
        std::lock_guard<std::mutex> lock(m);
        return data.empty();
    }
};
\end{lstlisting}

\subsection{Q1: Thread Safe Queue with notification}
\begin{itemize}
    \item Analysis from stack also applies
    \item \verb|wait_and_pop()| provides a solution to waiting for an entry
    \item If a thread throws an exception after it is woken up, no other threads will be woken
    \item Can replace with a global notification, or notify one when an exception is thrown, or move shared pointer initialisation to push and store shared pointer instances instead
\end{itemize}

\begin{lstlisting}[language=C++]
template<typename T>
class threadsafe_queue {
private:
    mutable std::mutex mut;
    std::queue<T> data_queue;
    std::condition_variable data_cond;
public:
    threadsafe_queue() {}
    void push(T new_value) {
        std::lock_guard<std::mutex> lk(mut);
        data_queue.push(std::move(new_value));
        data_cond.notify_one();
    }
    void wait_and_pop(T& value) {
        std::unique_lock<std::mutex> lk(mut);
        data_cond.wait(lk, [this]{ return !data_queue.empty(); });
        value = std::move(data_queue.front());
        data_queue.pop();
    }
    std::shared_ptr<T> wait_and_pop() {
        std::unique_lock<std::mutex> lk(mut);
        data_cond.wait(lk, [this]{ return !data_queue.empty(); });
        std::shared_ptr<T> res(
            std::make_shared<T>(std::move(data_queue.front()))
        );
        data_queue.pop();
        return res;
    }
    bool try_pop(T& value) {
        std::lock_guard<std::mutex> lk(mut);
        if (data_queue.empty()) return false;
        value = std::move(data_queue.front());
        data_queue.pop();
        return true;
    }
    std::shared_ptr<T> try_pop() {
        std::lock_guard<std::mutex> lk(mut);
        if (data_queue.empty()) return std::shared_ptr<T>();
        std::shared_ptr<T> res(
            std::make_shared<T>(std::move(data_queue.front()))
        );
        data_queue.pop();
        return res;
    }
    bool empty() const {
        std::lock_guard<std::mutex> lk(mut);
        return data_queue.empty();
    }
};
\end{lstlisting}

\subsection{Q2: Push shared pointers in the queue}
\begin{itemize}
    \item Exception safe
    \item Allocaiton of new instance can be done outside of the lock for better performance
    \item Concurrency is still limited by global mutex
\end{itemize}

\begin{lstlisting}[language=C++]
template<typename T>
class threadsafe_queue {
private:
    mutable std::mutex mut;
    std::queue<std::shared_ptr<T>> data_queue;
    std::condition_variable data_cond;
public:
    threadsafe_queue() {}
    void wait_and_pop(T& value) {
        std::unique_lock<std::mutex> lk(mut);
        data_cond.wait(lk, [this]{ return !data_queue.empty(); });
        value = std::move(*data_queue.front());
        data_queue.pop();
    }
    bool try_pop(T& value) {
        std::lock_guard<std::mutex> lk(mut);
        if (data_queue.empty()) return false;
        value = std::move(*data_queue.front());
        data_queue.pop();
        return true;
    }
    std::shared_ptr<T> wait_and_pop() {
        std::unique_lock<std::mutex> lk(mut);
        data_cond.wait(lk, [this]{ return !data_queue.empty(); });
        std::shared_ptr<T> res = data_queue.front();
        data_queue.pop();
        return res;
    }
    std::shared_ptr<T> try_pop() {
        std::lock_guard<std::mutex> lk(mut);
        if (data_queue.empty()) return std::shared_ptr<T>();
        std::shared_ptr<T> res = data_queue.front();
        data_queue.pop();
        return res;
    }
    void push(T new_value) {
        std::shared_ptr<T> data(
            std::make_shared<T>(std::move(new_value))
        );
        std::lock_guard<std::mutex> lk(mut);
        data_queue.push(data);
        data_cond.notify_one();
    }
    bool empty() const {
        std::lock_guard<std::mutex> lk(mut);
        return data_queue.empty();
    }
};
\end{lstlisting}

\subsection{Q3: Single threaded queue}
\begin{itemize}
    \item Uses unique pointers to manage nodes so that they get deleted by RAII
    \item Pushing can modify both the front and back
    \item Pushing and popping access the next pointer of a node which requires serialisation
\end{itemize}

\begin{lstlisting}[language=C++]
template<typename T>
class queue {
private:
    struct node {
        T data;
        std::unique_ptr<node> next;
        node(T data_):
            data(std::move(data_))
        {}
    };
    std::unique_ptr<node> front;
    node* back;
public:
    queue(): back(nullptr) {}
    queue(const queue& other)=delete;
    queue& operator=(const queue& other)=delete;
    std::shared_ptr<T> try_pop() {
        if (!front) {
            return std::shared_ptr<T>();
        }
        std::shared_ptr<T> const res(
            std::make_shared<T>(std::move(front->data))
        );
        std::unique_ptr<node> const old_front = std::move(front);
        front = std::move(old_front->next);
        if (!front) back = nullptr;
        return res;
    }
    void push(T new_value) {
        std::unique_ptr<node> p(new node(std::move(new_value)));
        node* const new_back = p.get();
        if (back) back->next = std::move(p);
        else front = std::move(p);
        back = new_back;
    }
};
\end{lstlisting}

\subsection{Q4: Single threaded queue, Enabling concurrency by separating data}
\begin{itemize}
    \item Pre-allocate dummy node with no data, so that there is always at least one node in the queue separating the front and back nodes
    \item Extra level of indirection needed to store data by pointers
\end{itemize}

\begin{lstlisting}[language=C++]
template<typename T>
class queue {
private:
    struct node {
        std::shared_ptr<T> data;
        std::unique_ptr<node> next;
    };
    std::unique_ptr<node> front;
    node* back;
public:
    queue(): front(new node), back(front.get()) {}
    queue(const queue& other)=delete;
    queue& operator=(const queue& other)=delete;
    std::shared_ptr<T> try_pop() {
        if (front.get() == back) {
            return std::shared_ptr<T>();
        }
        std::shared_ptr<T> const res(front->data);
        std::unique_ptr<node> old_front = std::move(front);
        front = std::move(old_front->next);
        return res;
    }
    void push(T new_value) {
        std::shared_ptr<T> new_data(
            std::make_shared<T>(std::move(new_value))
        );
        std::unique_ptr<node> p(new node);
        back->data = new_data;
        node* const new_back = p.get();
        back->next = std::move(p);
        back = new_back;
    }
};
\end{lstlisting}

\subsection{Q5: Fine-grained Thread-safe Queue}
Invariants Held:
\begin{itemize}
    \item \verb|back->next == nullptr|
    \item \verb|back->data == nullptr|
    \item \verb|front == back| implies an empty list
    \item A single element list has \verb|front->next == back|
    \item For each node \verb|x| where \verb|x != back|, \verb|x->data| points to an instance of \verb|T| and \verb|x->next| points to the next node in the list; \verb|x->next == back| implies \verb|x| is the last node in the list
    \item Following the \verb|next| nodes from \verb|front| will eventually yield \verb|back|
\end{itemize}

\begin{lstlisting}[language=C++]
template<typename T>
class threadsafe_queue {
private:
    struct node {
        std::shared_ptr<T> data;
        std::unique_ptr<node> next;
    };
    std::mutex front_mutex;
    std::unique_ptr<node> front;
    std::mutex back_mutex;
    node* back;
    node* get_back() {
        std::lock_guard<std::mutex> back_lock(back_mutex);
        return back;
    }
    std::unique_ptr<node> pop_front() {
        std::lock_guard<std::mutex> front_lock(front_mutex);
        if (front.get() == get_back()) {
            return nullptr;
        }
        std::unique_ptr<node> old_front = std::move(front);
        front = std::move(old_front->next);
        return old_front;
    }
public:
    threadsafe_queue(): front(new node), back(front.get()) {}
    threadsafe_queue(const queue& other)=delete;
    threadsafe_queue& operator=(const queue& other)=delete;
    std::shared_ptr<T> try_pop() {
        std::unique_ptr<node> old_front = pop_front();
        return old_front ? old_front->data : std::shared_ptr<T>();
    }
    void push(T new_value) {
        std::shared_ptr<T> new_data(
            std::make_shared<T>(std::move(new_value))
        );
        std::unique_ptr<node> p(new node);
        node* const new_back = p.get();
        std::lock_guard<std::mutex> back_lock(back_mutex);
        back->data = new_data;
        back->next = std::move(p);
        back = new_back;
    }
};
\end{lstlisting}

\subsection{Q6: Even more Fine-grained Lock-based Queue}
\begin{itemize}
    \item Per-node mutexes provide synchronises with relationship between push and pop threads
    \item Additional scope introduced when trying to pop avoids UAF
\end{itemize}

\begin{lstlisting}[language=C++]
class FineQueue {
    struct Node {
        std::mutex mut;
        Node *next = nullptr;
        Data data{};
    };
    std::mutex mut_back;
    Node *back;
    std::mutex mut_front;
    Node *front;
public:
    FineQueue() : mut_back{},
                  back{new Node{}},
                  mut_front{},
                  front{back} {}
    ~FineQueue() {
        while (front != nullptr) {
            Node *next = front->next;
            delete front;
            front = next;
        }
    }
    void push(Data data) {
        Node *new_node = new Node{};
        std::scoped_lock lock{mut_back};
        std::scoped_lock node_lock{back->mut};
        back->data = data;
        back->next = new_node;
        back = new_node;
    }
    std::optional<Data> try_pop() {
        Node *old_node;
        {
            std::scoped_lock lock{mut_front};
            old_node = front;
            std::scoped_lock node_lock{old_node->mut};
            if (old_node->next == nullptr) {
                return std::nullopt;
            }
            front = front->next;
        }
        Data data = old_node->data;
        delete old_node;
        return data;
    }
};
\end{lstlisting}

\section{Testing and Debugging}

\subsection{Using Valgrind}
\begin{lstlisting}
valgrind --tool=memcheck <prog>
\end{lstlisting}

\subsection{Using ASan}
\begin{lstlisting}
clang++ -O1 -fsanitize=address a.cc && ./a.out
\end{lstlisting}

\subsection{How ASan works}
\begin{lstlisting}
// X is good byte, O is bad byte, mapped to shadow value
XXXXXXXX  0
XXXXXXXO  7
XXXXXXOO  6
XXXXXOOO  5
XXXXOOOO  4
XXXOOOOO  3
XXOOOOOO  2
XOOOOOOO  1
OOOOOOOO -1

*a=...;
char *shadow = addr >> 3; // find corresponding byte
if (*shadow // all safe has value zero, fails this check
&& *shadow <= ((a&7)+N-1)) { // find offset within block and last accessed byte
    ReportError(a);
}
*a=...;
\end{lstlisting}

\subsection{Using TSan}
\begin{lstlisting}
g++ a.cpp -fsanitize=thread -fPIE -pie
TSAN_OPTIONS="history-size=7 force_seq_cst_atomics=1" ./myprogram
\end{lstlisting}

\subsection{How TSan works}
\begin{lstlisting}
Shadow = 4 * (Addr & kMask); // shadow address hash function

TID: 16 bits (TID, thread ID)
Epoc: 42 bits (Epoch, scalar clock)
Pos: 5 bits (Position in 8-byte word)
IsW: 1 bit (IsWrite)
\end{lstlisting}

\section{Midterms Tutorial Stuff}

\subsection{C++ Concurrent Queue with Semaphore}

\begin{lstlisting}[language=C++]
struct Job {
  int id;
  int count;
};

class JobQueue {
  std::queue<Job> jobs;
  std::mutex mut;
  std::counting_semaphore<> count;

public:
  JobQueue() : jobs{}, mut{}, count{0} {}

  void enqueue(Job job) {
    std::unique_lock lock{mut};
    jobs.push(job);
    count.release(); // increment semaphore counter
  }

  std::optional<Job> try_dequeue() {
    if (count.try_acquire()) {
      // enter critical section and remove a job
      std::unique_lock lock{mut};
      Job job = jobs.front();
      jobs.pop();
      return job;
    } else {
      // queue is empty
      return std::nullopt;
    }
  }

  Job dequeue() {
    count.acquire(); // decrement semaphore counter
    std::unique_lock lock{mut};
    Job job = jobs.front();
    jobs.pop();
    return job;
  }
};
\end{lstlisting}

\subsection{C++ Concurrent Queue with Monitor}

\begin{lstlisting}[language=C++]
struct Job {
    int id;
    int count;
}
class JobQueue {
    std::queue<Job> jobs;
    std::mutex mut;
    std::condition_variable cond;
public:
    JobQueue() : jobs{}, mut{}, cond{} {}
    void enqueue(Job job) {
        {
            std::unique_lock lock{mut};
            jobs.push(job);
        }
        cond.notify_one();
    }
    std::optional<Job> try_dequeue() {
        std::unique_lock lock{mut};
        if (jobs.empty()) {
            return std::nullopt;
        } else {
            Job job = jobs.front(); jobs.pop();
            return job;
        }
    }
    Job dequeue() {
        std::unique_lock lock{mut};
        cond.wait(lock, [this]() { return !jobs.empty(); });
        // same as
        // while (jobs.empty()) {
        //     cond.wait(lock);
        // }
        Job job = jobs.front();
        jobs.pop();
        return job;
    }
};
\end{lstlisting}

\subsection{Atomic Puzzle Tips}
\begin{itemize}
    \item Find possible modification orders; they must obey write-write coherence.
    \item Decide the possible outcomes of reads; they must obey read-write, write-read, write-write coherence.
    \item Deduce the sw relationships between threads; the load must successfully read the store (or value in release sequence) and the two operations must be suitably tagged.
    \item Deduce the hb relationships.
    \item For non-atomic variables, any two operations involving at least one write must be synchronised; otherwise, there is a data race.
\end{itemize}

\subsection{Smart Pointers}
\begin{itemize}
    \item Unique pointer implements RAII with unique ownership. 
    \item \verb|std::shared_ptr| stores \verb|ref_count| (thread safe) and the underlying object (not thread safe), protecting shared resources.
    \item Weak pointer acts as a viewer and helps to prevent circular dependencies since it does not contribute to ownership.
    \item The underlying object is deleted when \verb|shared_count| hits 0.
    \item The control block is deleted when \verb|weak_count| hits 0.
    \item If \verb|shared_count| is above 0, it counts towards \verb|weak_count| as well.
\end{itemize}

\begin{lstlisting}[language=C++]
std::unique_ptr<int> foo = std::make_unique<int>(0);
foo.get(); // *foo

std::shared_ptr<int> x = std::make_shared<int>(100);
{
    auto a = x;
    auto b = x;
    cout << x.use_count();
}

// Shared Pointer Usage

// Stack (Local Variables)
shared_ptr { ctrl_blk* }
weak_ptr   { ctrl_blk* }

// Heap (The Control Block)
ctrl_blk {
    std::atomic<size_t> shared_count; // Controls object life
    std::atomic<size_t> weak_count;   // Controls control block life
    T* ptr_to_obj;                    // The actual data
}
\end{lstlisting}

\subsection{C++ Shared Pointer}
\begin{lstlisting}[language=C++]
template <typename T>
struct SharedPtr {
private:
    std::atomic<size_t>* m_count;
    T* m_ptr;
public:
    SharedPtr(T* ptr) : m_count(new std::atomic<size_t>(1)), m_ptr(ptr) { }
    SharedPtr(const SharedPtr& other) : m_count(other.m_count),                                   m_ptr(other.m_ptr) {
        m_count->fetch_add(1, std::memory_order::relaxed);
    }
    ~SharedPtr() {
        size_t old_count = m_count->fetch_sub(1, std::memory_order::acq_rel);
        if(old_count == 1) {
            delete m_ptr;
            delete m_count;
        }
    }
    T* get() { return m_ptr; }
    const T* get() const { return m_ptr; }
};
\end{lstlisting}

\subsection{Compare and Swap}
\begin{lstlisting}[language=C++]
int v = x.load();
while (true) {
    // could spuriously fail
    if (x.compare_and_exchange_weak(v, v+1)) {
        return v;
    }
}
// or
x.compare_and_exchange_strong(v, v+1);
// similar to
x.fetch_add(1);
\end{lstlisting}

\subsection{C++ Lock Free Stack}

\begin{lstlisting}[language=C++]
void push(Job job) {
    Node* newTop = new Node{job, nullptr};
    Node* oldTop = top.load();
    while (true) {
        newTop->next.store(oldTop);
        if (top.compare_exchange_weak(oldTop, newTop)) {
            break;
        }
    }
}
std::optional<Job> try_pop() {
    Node* oldTop = top.load();
    while (true) {
        if (oldTop == nullptr) return std::nullopt;
        Node* newTop = oldTop->next.load();
        if (top.compare_exchange_weak(oldTop, newTop)) {
            Job job = oldTop->job;
            delete oldTop;
            return job;
        }
    }
}
\end{lstlisting}

\subsection{C++ Lock Free Queue + Recycling Centre}
\begin{itemize}
    \item Unbounded lock free queue with \verb|push| and non-blocking \verb|try_pop|
    \item Avoids ABA problem using generation counter
    \item Avoids UAF by reusing nodes instead of deallocating them
    \item However, does not provide serialisation
\end{itemize}

\begin{lstlisting}[language=C++]
struct Job
{
    int id;
    int data;
};
class JobQueue {
    using stdmo = std::memory_order;

    // a node is a dummy node if its next pointer is set to QUEUE_END
    // use the next ptr to establish the synchronizes-with relationship
    // next is in charge of releasing job
    struct Node {
        std::atomic<Node*> next = QUEUE_END;
        Job job;
    };

    // sentinel values
    static inline Node* const QUEUE_END = nullptr;
    static inline Node* const STACK_END = QUEUE_END + 1;

    struct alignas(16) GenNodePtr {
        Node* node;
        uintptr_t gen;
    };

    static_assert(std::atomic<GenNodePtr>::is_always_lock_free);

    alignas(64) std::atomic<Node*> m_queue_back;              // producer end
    alignas(64) std::atomic<GenNodePtr> m_queue_front;        // consumer end
    alignas(64) std::atomic<GenNodePtr> m_recycled_stack_top; // recycled node stack

public:
    // queue starts with a dummy node
    JobQueue()
        : m_queue_back(new Node())
        , m_queue_front(GenNodePtr { m_queue_back.load(stdmo::relaxed), 1 })
        , m_recycled_stack_top(GenNodePtr { STACK_END, 1 }) {}

    ~JobQueue() {
        // assumes no other threads are accessing the job queue
        Node* cur_node = m_queue_front.load(stdmo::relaxed).node;
        while(cur_node != QUEUE_END) {
            Node* next = cur_node->next;
            delete cur_node;
            cur_node = next;
        }

        // clean up the recycled nodes
        cur_node = m_recycled_stack_top.load(stdmo::relaxed).node;
        while(cur_node != STACK_END) {
            Node* next = cur_node->next;
            delete cur_node;
            cur_node = next;
        }
    }

    // either get a node from the recycling stack if we have some
    // or allocate a new one if we dont
    Node* get_recycled_node_or_allocate_new() {
        GenNodePtr old_stack_top = m_recycled_stack_top.load(stdmo::relaxed);
        while(true) {
            if(old_stack_top.node == STACK_END) {
                return new Node();
            }

            // use acquire, synchronise with the release-store of
            // node->next in add_node_to_recycling_stack
            Node* cur_stack_next = old_stack_top.node->next.load(stdmo::acquire);

            GenNodePtr new_stack_top { cur_stack_next, old_stack_top.gen + 1 };

            if(m_recycled_stack_top.compare_exchange_weak(
                   old_stack_top,
                   new_stack_top,
                   stdmo::relaxed)) {
                return old_stack_top.node;
            }
        }
    }

    // put node in recycling centre
    void add_node_to_recycling_stack(Node* node) {
        GenNodePtr old_stack_top = m_recycled_stack_top.load(stdmo::relaxed);
        while(true) {
            // use release, synchronise with the acquire-load of
            // node->next in get_recycled_node_or_allocate_new
            node->next.store(old_stack_top.node, stdmo::release);

            GenNodePtr new_stack_top { node, old_stack_top.gen + 1 };

            if(m_recycled_stack_top.compare_exchange_weak(
                   old_stack_top,
                   new_stack_top,
                   stdmo::relaxed)) {
                break;
            }
        }
    }
public:
    void push(Job job) {
        Node* new_dummy = get_recycled_node_or_allocate_new();
        new_dummy->next.store(QUEUE_END, stdmo::relaxed);

        Node* work_node = m_queue_back.exchange(new_dummy, stdmo::acq_rel);
        work_node->job = job;
        work_node->next.store(new_dummy, stdmo::release);
    }

    std::optional<Job> try_pop() {
        // splice node from the front of queue, but only if its not dummy
        // successfully splicing a node establishes global order of pops

        GenNodePtr old_front = m_queue_front.load(stdmo::acquire);
        while(true) {
            // get node out from the GenNodePtr
            Node* old_front_next = old_front.node->next.load(stdmo::acquire);
            if(old_front_next == QUEUE_END) {
                return std::nullopt;
            }

            // note that the generation is strictly increasing
            GenNodePtr new_front { old_front_next, old_front.gen + 1 };

            if(m_queue_front.compare_exchange_weak(old_front,
                   new_front, stdmo::acq_rel)) {
                // node now belongs to us
                break;
            }
        }

        Job job = old_front.node->job;
        add_node_to_recycling_stack(old_front.node);
        return job;
    }
};
\end{lstlisting}

\section{Go Basics}

\subsection{Fundamentals}
\begin{lstlisting}[language=Golang]
world := "World"
var exclamation string = "!" // name before type
defer func(){}() // RAII style, FILO

if res, err := something(); err != nil {

}

s := make([]int, 0, 5) // size 0, cap 5
s = append(s, 0) // append 0 to slice

m1 := make(map[string]int)
m2 := make[string]int{"A": 65}
val, ok := m1[key1]
// iteration order is random

// set trick
s := map[string]struct{}{}
_, exist := s["A"]

for i := 0; i < 5; i++ {}
for {
    if ok := check(); !ok { break }
}
for check() {}
for i, elem := range[]int{1, 2, 3} {} // use _ if unwanted
for k, v := range map[int]int{} {}
for elem := range elemCh {}

switch num {
case 1: smth()
case 2: fallthrough
default: default()
}

type Duck interface { Walk() }
type SomeBird struct { wing Wing }
func (b *SomeBird) Walk() {} // tied to an instance, by reference
// SomeBird fulfils Duck even if not explictly linked
\end{lstlisting}

\subsection{Goroutines}
\begin{lstlisting}[language=Golang]
func hello() {
    fmt.Println("hello")
}
func main() {
    go hello()
    go func() { fmt.Println("bye") }() // anonymous function
}
\end{lstlisting}

\subsection{Synchronising Goroutines}
Use \texttt{WaitGroup}.
\begin{lstlisting}[language=Golang]
func main() {
    var wg sync.WaitGroup
    for _, x := range []string{"a", "b", "c"} {
        wg.Add(1)
        // closure over x
        go func() {
            defer wg.Done() // RAII style
            fmt.Println(x)
        }
        // copy x
        go func(x string) {
            defer wg.Done()
            fmt.Println(x)
        }
    }
    wg.Wait()
}
\end{lstlisting}

\subsection{Channels}
\begin{lstlisting}[language=Golang]
// bidirectional channels
var x chan interface{}
x = make(chan interface{})

// unidirectional channels
var rx <-chan interface{}
var tx chan<- interface{}
x := make(chan interface{})
rx = x
tx = x

// buffered channel
c := make(chan interface{}, 4)

// using channels
s := make(chan string)
go func() {
    s <- "hi"
}()
t, ok := <-s
// ok indicates whether read was a write or default value
// reading from closed channel is always allowed
\end{lstlisting}

\subsection{Synchronisation with Channels}
\begin{lstlisting}[language=Golang]
// ranging over channel
c := make(chan int)
go func() {
    defer close(c)
    for i := 1; i <= 5; i++ {
        c <- i
    }
}()
for x := range c {
    fmt.Println(x)
}

c := make(chan interface{})
var wg sync.WaitGroup
for i := 0; i < 5; i++ {
    wg.Add(1)
    go func(i int) {
        defer wg.Done()
        <-c
    }
}
close(c)
wg.Wait()
\end{lstlisting}

\subsection{\texttt{select}}
\begin{lstlisting}[language=Golang]
// pseudorandom uniform selection
// no fall through
c1 := make(chan interface{}); close(c1)
c2 := make(chan interface{}); close(c2)
var x1, x2 int
for i := 1000; i >= 0; i-- {
    select {
        case <- c1: x1++
        case <- c2: x2++
        default:
    }
}

// timeout channel implementation
select {
    case <-c:
    case <- time.After(1 * time.Second):
        fmt.Println("Timed out")
}

// for-select
loop:
for {
    select {
        case <- done: break loop
        default:
    }
    // do work
}
\end{lstlisting}

\subsection{Synchronisation Constructs}
\begin{lstlisting}[language=Golang]
// RWMutex
var mu sync.RWMutex
mu.RLock() // multiple readers can hold
defer mu.RUnlock()

// similar to Mutex
mu.Lock() // only one writer can hold
defer mu.Unlock()

// Once
var once sync.Once
// excutes only once
once.Do(func() {
    // something
})

// Cond
mu := &sync.Mutex{}
cond := sync.NewCond(mu)
ready := false
// waiter
go func() {
    mu.Lock()
    for !ready {
        cond.Wait()
    }
    fmt.Println("Ready!")
    mu.Unlock()
}()
// In the signaller
mu.Lock()
ready = true
cond.Signal() // wake one
cond.Broadcast() // wake all
mu.Unlock()

// Pool
var bufferPool = sync.Pool{
    New: func() interface{} {
        return make([]byte, 1024)
    },
}
func process() {
    buf := bufferPool.Get().([]byte)
    defer bufferPool.Put(buf)
    // use buffer
}
\end{lstlisting}

\section{Go Patterns}

\subsection{Lexical Confinement}
\begin{lstlisting}[language=Golang]
// only expose reading end of channel to others
owner := func() <-chan int {
    res := make(chan int, 5)
    go func() {
        defer close(res)
        for i := 0; i <= 5; i++ {
            res <- i
        }
    }()
    return res
}   
consumer := func(res <-chan int) {
    for res := range res {
        fmt.Println(res)
    }
}
res := owner()
consumer(res)

printData := func(data []byte) {
    // ...
}
// expose only slices of arrays
go printData(data[:3])
go printData(data[3:])
\end{lstlisting}

\subsection{For-Select Loop}
\begin{lstlisting}[language=Golang]
for {
    select {
        case num := <-q:
            sum += num
        case <-done:
            sumCh <- sum
            return
    }
}
\end{lstlisting}

\subsection{Preventing Leaking}
\begin{lstlisting}[language=Golang]
// returns status channel
doWork := func(
    done <-chan interface{},
    strings <- chan string,
) <-chan interface{} {
    terminated := make(chan interface{})
    go func() {
        defer close(terminated)
        for {
            select {
                case s := <-strings: fmt.Println(s)
                case <-done: return
            }
        }
    }
}

// done channel
done := make(chan interface{})
terminated := doWork(done, nil)
go func() {
    time.Sleep(1 * time.Second)
    close(done)
}
<-terminated
\end{lstlisting}

\subsection{Error Handling}
\begin{lstlisting}[language=Golang]
type Result struct {
    Error error
    Response *http.Response
}
checkStatus := func(done <-chan interface{}, urls ...string) <-chan Result {
    results := make(chan Result)
    go func() {
        defer close(results)
        for _, url := range urls {
            var result Result
            resp, err := http.Get(url)
            result = Result{Error: err, Response: resp}
            select {
                case <-done: return
                case results <- result:
            }
        }
        return results
    }
}
done := make(chan interface{})
defer close(done)
urls := []string{"", ""}
for res := range checkStatus(done, urls...) {
    if res.Error != nil {
        continue
    }
}
\end{lstlisting}

\subsection{Pipeline}
\begin{lstlisting}[language=Golang]
generator := func(done <-chan interface{}, integers ...int) <-chan int {
    intStream := make(chan int)
    go func() {
        defer close(intStream)
        for _, i := range integers {
            select {
                case <-done: return
                case intStream <- i:
            }
        }
    }()
    return intStream
}
multiply := func(
    done <-chan interface{},
    intStream <-chan int,
    multiplier int,
) <-chan int {
    multipliedStream := make(chan int)
    go func() {
        defer close(multipliedStream)
        for i := range intStream {
            select {
                case <-done: return
                case multipliedStream <- i * multiplier:
                // add is the same but with + here
            }
        }
    }()
    return multipliedStream
}

done := make(chan interface{})
defer close(done)
intStream := generator(done, 1, 2, 3, 4)
// * 2 + 1 * 2
pipeline := multiply(done, add(done, multiply(done, intStream, 2), 1), 2)
for v := range pipeline {
    fmt.Println(v)
}
\end{lstlisting}

\subsection{Fan-out, Fan-in}
\begin{lstlisting}[language=Golang]
fanIn := func(
    done <-chan interface{},
    channels ...<-chan interface{}
) <-chan interface {
    var wg sync.WaitGroup
    multiplexedStream := make(chan interface{})
    multiplex := func(c <-chan interface{}) {
        defer wg.Done()
        for i := range c {
            select {
                case <-done: return
                case multiplexedStream <- i:
            }
        }
    }
    wg.Add(len(channels))
    for _, c := range channels {
        go multiplex(c)
    }
    go func() {
        wg.Wait()
        close(multiplexedStream)
    }()
    return multiplexedStream
}

done := make(chan interface{})
defer close(done)
start := time.Now()
rand := func() interface{} { return rand.Intn(676767) }
randIntStream := toInt(done, repeatFn(done, rand))
numFinders := runtime.NumCPU() // number of OS threads
finders := make([]<-chan interface{}, numFinders)
for i := 0; i < numFinders; i++ {
    finders[i] = primeFinder(done, randIntStream)
}
// limit output
for prime := range take(done, fanIn(done, finders...), 10) {
    fmt.Println(prime)
}
\end{lstlisting}

\subsection{Load Balancer}
\begin{lstlisting}[language=Golang]
// requests sent to balancer
type Request struct {
    fn func() int // operation to perform
    c chan int // channel for result
}

// requester
func requester(work chan<- Request) {
    c := make(chan int)
    for {
        work <- Request{ workFn, c }
        result := <-c
        process(result)
    }
}

// worker
type Worker struct {
    requests chan Request // buffered work channel
    pending int // count of tasks
    index int // index in heap
}

func (w *Worker) work(done chan *Worker) {
    for {
        req := <-w.requests
        req.c <- req.fn()
        done <- w
    }
}

// balancer
type Pool []*Worker
type Balancer struct {
    pool Pool
    done chan *Worker
}

func (b *Balancer) balance(work chan Request) {
    for {
        select {
            case req := <-work: b.dispatch(req)
            case w := <-b.done: b.completed(w)
        }
    }
}

// worker pool
func (p Pool) Less(i, j int) bool {
    return p[i].pending < p[j].pending
}

// dispatcher
func (b *Balancer) dispatch(req Request) {
    w := heap.Pop(&b.pool).(*Worker)
    w.requests <- req
    w.pending++
    heap.Push(&b.pool, w)
}

// completed
func (b *Balancer) completed(w *Worker) {
    w.pending--
    heap.Remove(&b.pool, w.index)
    heap.Push(&b.pool, w)
}
\end{lstlisting}

\section{Classical Synchronisation Problems}

\subsection{Producer Consumer}
\begin{lstlisting}
// producer
Produce Item;
wait(notFull);
wait(mutex);
buffer[in] = item;
in = (in+1) % K;
count++;
signal(mutex);
signal(notEmpty);

// consumer
wait(notEmpty);
wait(mutex);
item = buffer[out];
out = (out+1) % K;
count--;
signal(mutex);
signal(notFull);
Consume Item;
\end{lstlisting}

\subsection{Readers Writers}
\begin{lstlisting}
// lightswitch
counter = 0
mutex = Semaphore(1)
roomEmpty = Semaphore(1)
lock(semaphore):
    mutex.wait()
    counter += 1
    if counter == 1:
        semaphore.wait()
    mutex.signal()
unlock(semaphore):
    mutex.wait()
    counter -= 1
    if counter == 0:
        semaphore.signal()
    mutex.signal()

// writer
turnstile.wait()
roomEmpty.wait()
// CS
turnstile.signal()
roomEmpty.signal()

// reader
turnstile.wait()
turnstile.signal()
readLightswitch.lock(roomEmpty)
// CS
readLightswitch.unlock(roomEmpty)
\end{lstlisting}

\subsection{\texttt{std::barrier}}
\begin{lstlisting}[language=C++]
std::barrier sync_point(std::ssize(x), complete);
auto work = [&](std::string y) {
    // work
    sync_point.arrive_and_wait();
    // work
    sync_point.arrive_and_wait();
}
std::vector<std::thread> threads;
for (auto const& y : x) threads.emplace_back(work, y);
for (auto& thread : threads) thread.join();
\end{lstlisting}

\subsection{Barrier Implementation}

\subsection{C++}
\begin{lstlisting}[language=C++]
struct Barrier {
    std::ptrdiff_t expected;
    std::ptrdiff_t count;
    std::mutex mut;
    std::counting_semaphore<> turnstile;
    std::counting_semaphore<> turnstile2;
    Barrier(std::ptrdiff_t expected)
        : expected{expected}, count{0}, mut{}, turnstile{0}, turnstile2{0} {}
    void arrive_and_wait() {
        {
            std::scoped_lock lock{mut};
            count++;
            if (count == expected) {
                turnstile2.acquire(); // close waiter turnstile
                turnstile.release(expected); // open turnstile into CS
            }
        }
        turnstile.acquire();
        {
            std::scoped_lock lock{mut};
            count--;
            if (count == 0) {
                turnstile.acquire(); // close turnstile to reset barrier
                turnstile2.release(expected); // open waiter turnstile
            }
        }
        turnstile2.acquire();
    }
};
\end{lstlisting}

\subsection{Golang}
\begin{lstlisting}[language=Golang]
type Barrier struct {
    wg sync.WaitGroup
    wg2 sync.WaitGroup
}

func (b *Barrier) Init(expected int) {
    b.wg.Add(expected)
    b.wg2.Add(expected)
}

func (b *Barrier) Wait() {
    b.wg.Done()
    b.wg.Wait()
    // all threads have called wait
    // reset barrier
    b.wg.Add(1)
    b.wg2.Done()
    // wait for barrier to fully reset
    b.wg2.Wait()
    // barrier reset
    b.wg2.Add(1)
}
\end{lstlisting}

\subsection{Dining Philosophers}

\subsubsection{C++}
\begin{lstlisting}[language=C++]
// solution 1: basic solution
template<size_t NumP>
struct DiningTable {
    using ChpStick = std::mutex;
    ChpStick chpSticks[NumP];
    ChpStick& get_left_chpStick(size_t pid) { return chpSticks[pid]; }
    ChpStick& get_right_chpStick(size_t pid) { return chpSticks[(pid+1) % NumP]; }
    void eat(size_t pid, void (*eat_callback)(size_t pid)) {
        std::scoped_lock lock{get_left_chpStick(pid), get_right_chpStick(pid)};
        eat_callback(pid);
    }
};

// solution 2: limited eaters
template<size_t NumP>
struct DiningTable {
    using ChpStick = std::mutex;
    ChpStick chpSticks[NumP];
    std::counting_semaphore<> footman_sem{NumP-1};
    ChpStick& get_left_chpStick(size_t pid) { return chpSticks[pid]; }
    ChpStick& get_right_chpStick(size_t pid) { return chpSticks[(pid+1) % NumP]; }
    void eat(size_t pid, void (*eat_callback)(size_t pid)) {
        footman_sem.acquire();
        std::scoped_lock left_lock{get_left_chpStick(pid)};
        std::scoped_lock right_lock{get_right_chpStick(pid)};
        eat_callback(pid);
        footman_sem.release();
    }
};

// solution 3: fair semaphore
template<size_t NumP>
struct DiningTable {
    using ChpStick = FairMutex;
    ChpStick chpSticks[NumP];
    FairSemaphore footman_sem{NumP-1};
    ChpStick& get_left_chpStick(size_t pid) { return chpSticks[pid]; }
    ChpStick& get_right_chpStick(size_t pid) { return chpSticks[(pid+1) % NumP]; }
    void eat(size_t pid, void (*eat_callback)(size_t pid)) {
        footman_sem.acquire();
        std::scoped_lock left_lock{get_left_chpStick(pid)};
        std::scoped_lock right_lock{get_right_chpStick(pid)};
        eat_callback(pid);
        footman_sem.release();
    }
};

\end{lstlisting}

\subsubsection{Golang}
\begin{lstlisting}[language=Golang]
type ChpStick struct {}
type DiningTable struct {
    numPhilosophers int
    chpStickChs []chan ChpStick
}

func (t *DiningTable) Init(numPhilosophers int) {
    t.numPhilosophers = numPhilosophers
    t.chpStickChs = make([]chan ChpStick, 0, numPhilosophers)
    for i := 0; i < numPhilosophers; i++ {
        chpStick := make(chan ChpStick, 1)
        chpStick <- ChpStick{}
        t.chpStickChs = append(t.chpStickChs, chpStick)
    }
}

// solution 1: odd-even ring communication
func (t *DiningTable) Eat(pid int, eat_callback func(pid int)) {
    evenChpStickCh := t.evenChpStickCh(pid)
    oddChpStickCh := t.oddChpStickCh(pid)
    <-evenChpStickCh
    <-oddChpStickCh
    eat_callback(pid)
    evenChpStickCh <- ChpStick{}
    oddChpStickCh <- ChpStick{}
}

func (t *DiningTable) leftChpStickCh(pid int) chan ChpStick {
    return t.chpStickChs[pid]
}

func (t *DiningTable) rightChpStickCh(pid int) chan ChpStick {
    return t.chpStickChs[(pid+1) % t.numPhilosophers]
}

func (t *DiningTable) evenChpStick(pid int) chan ChpStick {
    if pid % 2 == 0 {
        return t.leftChpStickCh(pid)
    } else {
        return t.rightChpStickCh(pid)
    }
}

func (t *DiningTable) oddChpStick(pid int) chan ChpStick {
    if pid % 2 == 1 {
        return t.leftChpStickCh(pid)
    } else {
        return t.rightChpStickCh(pid)
    }
}

// solution 2: swap order of chopsticks to avoid deadlock
func (t * DiningTable) Eat(pid int, eat_callback func(pid int)) {
    evenChpStickCh := t.evenChpStickCh(pid)
    oddChpStickCh := t.oddChpStickCh(pid)
    
breakLock:
    for {
        <-evenChpStickCh
        select {
            case <-oddChpStickCh:
                break breakLock
            default:
                // swap
        }
        evenChStickCh <- ChpStick{}
        <-oddChpStickCh
        select {
            case <-evenChpStickCh:
                break breakLock
            default:
                /// swap
        }
    }
    eat_callback(pid)
    evenChpStickCh <- ChpStick{}
    oddChpStickCh <- ChpStick{}
}

// solution 3: limited eater using footman
type DiningTable struct {
    numPhilosophers int
    chpStickChs []chan ChpStick
    footman chan struct{}
}

func (t *DiningTable) Init(numPhilosophers int) {
    t.numPhilosophers = numPhilosophers
    t.chpStickChs = make([]chan ChpStick, 0, numPhilosophers)
    for i := 0; i < numPhilosophers; i++ {
        chpStick := make(chan ChpStick, 1)
        chpStick <- ChpStick{}
        t.chpStickChs = append(t.chpStickChs, chpStick)
    }
    t.footman = make(chan struct{}, numPhilosophers - 1)
    for i := 0; i < numPhilosophers - 1; i++ {
        t.footman <- struct{}{}
    }
}

func (t *DiningTable) Eat(pid int, eat_callback(func(pid int)) {
    leftChpStickCh := t.leftChpStickCh(pid)
    rightChpStickCh := t.rightChpStickCh(pid)
    <-t.footman
    select {
        case <-leftChpStickCh: <-rightChpStickCh
        case <-rightChpStickCh: <-leftChpStickCh
    }
    eat_callback(pid)
    leftChpStickCh <- ChpStick{}
    rightChpStickCh <- ChpStick{}
    t.footman <- struct{}{}
}
\end{lstlisting}

\subsubsection{Pseudocode}
\begin{lstlisting}
// Tanenbaum solution
int state[N];
Semaphore mutex = 1;
Semaphore s[N];

void philosopher(int i) {
    while (TRUE) {
        Think();
        takeChpSticks(i);
        Eat();
        putChpSticks(i);
    }
}

void takeChpSticks(i) {
    wait(mutex);
    state[i] = HUNGRY;
    safeToEat(i);
    signal(mutex);
    wait(s[i]);
}

void safeToEat(i) {
    if ((state[i] == HUNGRY) &&
        (state[LEFT] != EATING) &&
        (state[RIGHT] != EATING)) {
        state[i] = EATING;
        signal(s[i]);
    }
}

void putChpSticks(i) {
    wait(mutex);
    state[i] = THINKING;
    safeToEat(LEFT);
    safeToEat(RIGHT);
    signal(mutex);
}

// limited eater
// with N-1, guaranteed to have at least one success
void philosopher(int i) {
    while (TRUE) {
        Think();
        wait(seats);
        wait(chpStick[LEFT]);
        wait(chpStick[RIGHT]);
        Eat();
        signal(chpStick[LEFT]);
        signal(chpStick[RIGHT]);
        signal(seats);
    }
}
\end{lstlisting}

\subsection{Barbershop}
\subsection{Pseudocode}
\begin{lstlisting}
// pseudocode
customers = 0
mutex = Semaphore(1)
customer = Semaphore(0)
barber = Semaphore(0)
customerDone = Semaphore(0)
barberDone = Semaphore(0)
// customer
wait(mutex);
if (customers == n) {
    signal(mutex);
    exit();
}
customers += 1;
signal(mutex);
signal(customer);
wait(barber);
getHairCut();
signal(customeDone);
wait(barberDone);
wait(mutex);
customers -= 1;
signal(mutex);
// barber
while (TRUE) {
    wait(customer);
    signal(barber);
    cutHair();
    wait(customerDone);
    signal(barberDone);
}
\end{lstlisting}

\subsection{C++}
\begin{lstlisting}[language=C++]
template <size_t MaxCustomers>
struct Barbershop {
    size_t customers;
    std::mutex mut;
    std::counting_semaphore<> customer_sem;
    std::counting_semaphore<> barber_sem;
    std::counting_semaphore<> done_customer_sem;
    std::counting_semaphore<> done_barber_sem;
    void barber(void (*cutHair)()) {
        while (true) {
            customer_sem.acquire();
            barber_sem.release();
            cutHair();
            done_customer_sem.acquire();
            done_barber_sem.release();
        }
    }
    void customer(void (*balk)(), void (*getHairCut)()) {
        {
            std::scoped_lock lock{mut};
            if (customers == MaxCustomers + 1) {
                balk();
                return;
            }
            customers++;
        }
        customer_sem.release();
        barber_sem.acquire();
        getHairCut();
        done_customer_sem.release();
        done_barber_sem.acquire();
        {
            std::scoped_lock lock{mut};
            customers--;
        }
    }
};
\end{lstlisting}

\subsection{Golang}
\begin{lstlisting}[language=Golang]
// might not be FIFO
type Barbershop struct {
    chairs chan int
    barberAvailable chan struct{}
}

func (bs *Barbershop) Init(numChairs int) {
    bs.chairs = make(chan int, numChairs)
    bs.barberAvailable = make(chan struct{})
}

func (bs *Barbershop) Barber(cutHair func()) {
    for {
        bs.barberAvailable <- struct{}{}
        <-bs.chairs
        cutHair()
    }
}

func (bs *Barbershop) Customer(balk func(), getHairCut func()) {
    select {
        case bs.chairs <- 1:
            <-bs.barberAvailable
            getHairCut()
        default:
            balk()
    }
}
\end{lstlisting}

\section{Rust}

\subsection{Rust Basics}
\begin{lstlisting}[language=Rust]
fn add_one(i: i32) -> i32 { i+1 }

let x = 1; // immutable
let mut y = 1; // mutable
let z: i32 = 1;

// closures
|i| i + 2
|i: i32| -> i32 { i + 2 }

&str // char array like
String // string type buffer

// struct and enum literals
OrderDetails {
    price: 1,
    count: 1,
    id: 1
}

match order {
    Order::Buy(details) => smth(),
    Order::Sell(details) => smth(),
    Order::Cancel { id } => smth()
}

let v = 1;
match v {
    1 => ...,
    _ => ...
}

pub enum Result<T, E> {
    Ok(T),
    Err(E),
}

// optionals
pub enum Option<T> {
    None,
    Some(T),
}

let order_id = Some(3);
let Some(value) = order_id else {
    // panic
};

order_id.unwrap();

// ? operator unwraps 'Ok' or returns 'Err'
let content = std::fs::read_to_string("config.txt")?; 
let content = match std::fs::read_to_string("config.txt") { Ok(v) => v, Err(e) => return Err(e.into()) };

// vectors
let mut odd = Vec::new();
odd.push(1);
let numbers = (1..=100).collect::<Vec<u32>>();
let even: Vec<&u32> = numbers.iter().filter(|&i| i % 2 == 0).collect();
for n in &odd { } // borrow in loop
\end{lstlisting}

\subsection{Ownership and Borrowing}
\begin{lstlisting}[language=Rust]
// transferring ownership
fn main() {
    let mut book = Vec::new();
    book.push(...);
    book.push(...);
    publish(book);
    // ownership transferred, shallow copy of fields to other stack
    // original book cannot be used
    // publish(book); // compilation error, used of moved value
}
fn publish(book: Vec<String>) {
    // ...
} // book gets dropped here

// immutable borrow, shared reference
fn main() {
    let mut book = Vec::new();
    book.push(...);
    book.push(...);
    publish(&book);
    // reference borrowed
    publish(&book); // ok now
    // book is mutable here
}
fn publish(book: &Vec<String>) {
    // book is immutable here and in main
}

// mutable borrow
fn main() {
    let mut book = Vec::new();
    book.push(...);
    book.push(...);
    publish(&mut book);
    publish(&mut book);
}
// book doesn't need mut since itself doesn't need to be mutable
// the underlying reference is already mutable
fn publish(book: &mut Vec<String>) {
    // book is not usable in main at all
    book.push(...)
}

// return ownership
fn publish(book: Vec<String>) {
    book
}
\end{lstlisting}

\subsection{Threads}

\begin{lstlisting}[language=Rust]
let loc = thread::spawn(|| { // anonymous function syntax
    "world"
});
println!("Hello, {}!", loc.join().unwrap());

// move transfers ownership
// without move, rust will try to infer (not recommended)
let mut dst = Vec::new();
thread::spawn(move || {
    dst.push(3);
});
dst.push(3); // error, ownership was transferred
\end{lstlisting}

\subsection{Lifetimes}
\begin{lstlisting}[language=Rust]
// static lifetime, data does not borrow anything that could go out of scope
let s: &'static str = "hello";

// annotations
// x, y are valid for the lifetime of scope a
// returned reference will live at least as long as the shortest lifetime of x, y
fn longest<'a>(x: &'a str, y: &'a str) -> &'a str {
    if x.len() > y.len() {
        x
    } else {
        y
    }
}

// 2 different scopes
// z has different lifetime compared to x, y
fn longest_print_z<'a, 'b>(x: &'a str, y: &'a str, z: 'b str) -> &'a str {
    println!("{}", z);
    if x.len() > y.len() {
        x
    } else {
        y
    }
}

// reference counted type
// error, cannot be sent across threads
let v = Rc:: new(vec![1, 2, 3]);
let v2 = v.clone();
thread::spawn(move || {
    println!("{}", v.len());
});
another_fn(v2.clone);

// traits (interfaces)
trait Clone {
    fn clone(&self) -> Self;
}
impl<T: Clone> Clone for Vec<T> {
    fn clone(&self) -> Vec<T> {
        let mut v = Vec::new();
        for elem in self {
            v.push(elem.clone());
        }
        return v;
    }
}

// atomic reference counted
// only allows shared references
let v = Arc:: new(vec![1, 2, 3]);
let v2 = v.clone();
// move shared reference into thread
thread::spawn(move || {
    // v cannot be modified
    println!("{}", v.len());
});
another_fn(&v2);
\end{lstlisting}

\subsection{Rust Synchronisation}
\begin{lstlisting}[language=Rust]
// mutex
fn sync inc(counter: &Mutex<i32>) {
    let mut data MutexGuard<i32> = counter.lock();
    *data += 1;
}

// atomics
let number = AtomicUsize::new(10);
let prev = number.fetch_add(1, SeqCst);
assert_eq!(prev, 10);
let prev = number.swap(2, SeqCst);
assert_eq!(prev, 11);
assert_eq!(number.load(SeqCst), 2);

// MPSC FIFO queue
let (tx, rx) = mpsc::channel();
let tx2 = tx.clone();
thread::spawn(move || tx.send(5));
thread::spawn(move || tx2.send(4));
// order of printing is not specified
println!("{:?}", rx.recv));
println!("{:?}", rx.recv));
\end{lstlisting}

\subsection{Crossbeam}
\begin{lstlisting}[language=Rust]
// scoped threads
fn main() {
    let v = vec![1, 2, 3];
    std::thread::scope(|scope| {
        scope.spawn(|inner_scope| {
            // do stuff with v
        });
    }).unwrap();
    // v is back here
}

// producer consumer
fn main() {
    let (send_end, receive_end) = bounded(CHANNEL_CAPACITY);
    let mut threads = vec![];
    for _i in 0 .. NUM_THREADS {
        let send_end = send_end.clone();
        threads.push(
            thread::spawn(move || {
                for _k in 0 .. ITEMS_PER_THREAD {
                    let produced_value = produce_item();
                    send_end.send(produced_value).unwrap();
                }
            })
        )
    }
    for _j in 0 .. NUM_THREADS {
        let receive_end = receive_end.clone();
        threads.push(
            thread::spawn(move || {
                for _k in 0 .. ITEMS_PER_THREAD {
                    let to_consume = receive_end.recv().unwrap();
                    consume_item(to_consume);
                }
            })
        )
    }
    for t in threads {
        let _ = t.join();
    }
}

// exponential backoff
use crossbeam_utils::Backoff;
use std::sync_atomic::AtomicUsize;
use std::sync::atomic::Ordering::SeqCst;
fn fetch_mul(a: &AtomicUsize, b: usize) -> usize {
    let backoff = Backoff::new();
    loop {
        let val = a.load(SeqCst);
        if a.compare_and_swap(val, val.wrapping_mul(b), SeqCst) == val {
            return val;
        }
        backoff.spin();
    }
}

// rough implementation of spin wait
fn spin_wait(ready: &AtomicBool) {
    let backoff = Backoff::new();
    while !ready.load(SeqCst) {
        backoff.snooze();
    }
}
\end{lstlisting}

\subsection{Rayon}
\begin{lstlisting}[language=Rust]
// sequential maximum
fn main() {
    let vec = init_vector();
    let max = MIN;
    vec.iter().for_each(|n| {
        if *n > max {
            max = *n;
        }
    });
}
// with Rayon
fn main() {
    let vec = init_vector()
    let max = AtomicI64::new(MIN);
    vec.par_iter().for_each(|n| {
        loop {
            let old = max.load(Ordering::SeqCst);
            if *n <= old {
                break;
            }
            let returned = max.compare_and_swap(old, *n, Ordering::SeqCst);
            if returned == old {
                break;
            }
        }
    });
}
\end{lstlisting}

\subsection{Futures}
\begin{lstlisting}[language=Rust]
// simplified trait
trait Future {
    type Output;
    fn poll(&mut self, cx: &mut Context) -> Poll<Self::Output>;
}
enum Poll<T> {
    Ready(T),
    Pending,
}

// composition
let future = placeOnStove(meat)
    .then(|meat| cookOneSide(meat))
    .then(|meat| flip(meat))
    .then(|meat| cookOneSide(meat));

// clunky composition
fn addToInbox(email_id: u64, recipient_id: u64) -> impl Future<Output=Result<(), Error>> {
    loadMessage(email_id)
        .and_then(|message| get_recipient(message, recipient_id))
        .map(|(message, recipient)| recipient.verifyHasSpace(&message))
        .and_then(|(message, recipient)| recipient.addToInbox(message))
}

// async, returns a future
// await waits for a future and retrieves a value
async fn addToInbox(email_id: u64, recipient_id: u64) -> Result<(), Error> {
    let message = loadMessage(email_id).await?;
    let recipient = get_recipient(recipient_id).await?;
    recipient.verifyHasSpace(&message)?;
    recipient.addToInbox(message).await
}

// state management
enum AddToInboxState {
    NotYetStarted { email_id: u64, recipient_id: u64},
    WaitingLoadMessage { recipient_id: u64, state: LoadMessageFuture },
    WaitingGetRecipient { message: Message, state: GetRecipientFuture },
    WaitingAddToInbox { state: AddToInboxFuture },
    Completed { result: Result<(), Error> },
}

fn poll(*) {
    match self.state {
        NotYetStarted(email_id, recipient_id) => {
            let next_future = load_message(email_id);
            // switch to WaitingLoadMessage state
        },
        WaitingLoadMessage(email_id, recipient_id, state) => {
            match state.poll() {
                Ready(message) => {
                    let next_future = get_recipient(recipient_id);
                    // switch to WaitingGetRecipient state
                },
                Pending => return Pending,
            }
        },
        WaitingGetRecipient(message, recipient_id, state) => {
            match state.poll() {
                Ready(recipient) => {
                    recipient.verifyHasSpace(&message)?;
                    let next_future = recipient.addToInbox(message);
                    // switch to WaitingAddToInbox state
                },
                Pending => return Pending,
            }
        },
        // ...
    }
}
\end{lstlisting}

\section{TLA+}
\subsection{Basics}
\begin{lstlisting}
// AND: /\
// OR:  \/
// NOT: ~a

// predicate: expression returning boolean
operator(a, b, c) == (a /\ b) \/ (a /\ ~c)

// actions: tests
// tick means after
Next ==
    state = "hello"
    /\ state' = "goodbye"

// counting
Init ==
    x = 1
Next ==
    (x = 1 /\ x' = 2)
    \/ (x = 2 /\ x' = 3)
// alternatively, we can explicitly name the steps known as operators
// deadlocks at state 3
// staying at a state is not valid
Init == x = 1
Step1 == x = 1 /\ x' = 2
Step2 == x = 2 /\ x' = 3
Done == x = 3 /\ UNCHANGED x
// add new valid transitions
Next == Step1 \/ Step2 \/ Done \/ UNCHANGED x

// properties
[] (x > 0) // x always > 0
<> (x = 2) // eventually x = 2
<>[] (x = 3) // eventually x always = 3
(x = 2) ~> (x = 3) // x = 2 leads to x = 3
// stuttering causes these to fail

Spec = Init /\ [](Next \/ UNCHANGED x) // refactor specification
Spec = Init /\ [][Next]_x // syntax for next or stutter
Spec = Init /\ [][Next]_x /\ WF_x(Next) // fairness, if next step is available, eventually try it, ignoring stutter

// set theory
e \in S // inclusion
{1, 2, 3} // define set
{e \in S : p} // define set using predicate
\A e \in S : p // for all elements check predicate
\E x \in S : p // there exists element satisfying predicate
\end{lstlisting}

\subsection{Producer Consumer}
\begin{lstlisting}
// producer states and actions
CheckWritable ==
    producerState = "ready"
    /\ queueSize < MaxQueueSize
    /\ producerState' = "canWrite"
    /\ UNCHANGED queueSize

Write == 
    producerState = "canWrite"
    /\ producerState' = "ready"
    /\ queueSize' = queueSize + 1

ProducerAction == CheckWritable \/ Write

// consumer states and actions
CheckReadable ==
    consumerState = "ready"
    /\ queueSize > 0
    /\ consumerState' = "canRead"
    /\ UNCHANGED queueSize

Read == 
    consumerState = "canRead"
    /\ consumerState' = "ready"
    /\ queueSize' = queueSize - 1

ConsumerAction == CheckReadable \/ Read

// complete script
VARIABLES
    queueSize,
    producerState,
    consumerState

MaxQueueSize == 2

Init ==
    queueSize = 0
    /\ producerState = "ready"
    /\ consumerState = "ready"

Next == 
    ProducerAction
    \/ ConsumerAction
    \/ (UNCHANGED producerState
        /\ UNCHANGED consumerState
        /\ UNCHANGED queueSize)

// properties
AlwaysWithinBounds ==
    [] (queueSize >= 0
        /\ queueSize <= MaxQueueSize)

// multiple producers and consumers
// actually fails checks, since check and write a separate actions
CONSTANT producers, consumers
producers = {"p1", "p2"}
consumers = {"c1", "c2"}
VARIABLES queueSize, producerState, consumerState // maps / dictionaries
MaxQueueSize == 2
Init == 
    queueSize = 0
    /\ producerState = [p \in producers |-> "ready"]
    /\ consumerState = [c \in consumers |-> "ready"]

CheckWritable(p) == 
    producerState[p] = "ready"
    /\ queueSize < MaxQueueSize
    /\ producerState' =
        [producerState EXCEPT ![p] = "canWrite"]
    /\ UNCHANGED queueSize
    /\ UNCHANGED consumerState

Write(p) ==
    producerState[p] = "canWrite"
    /\ queueSize' = queueSize + 1
    /\ producerState' =
        [producerState EXCEPT ![p] = "ready"]
    /\ UNCHANGED consumerState

ProducerAction ==
    \E p \in producers : CheckWritable(p) \/ Write(p)

CheckReadable(c) == 
    consumerState[c] = "ready"
    /\ queueSize > 0
    /\ consumerState' =
        [consumerState EXCEPT ![c] = "canRead"]
    /\ UNCHANGED queueSize
    /\ UNCHANGED producerState

Read(c) ==
    consumerState[c] = "canRead"
    /\ queueSize' = queueSize - 1
    /\ consumerState' =
        [consumerState EXCEPT ![c] = "ready"]
    /\ UNCHANGED producerState

ConsumerAction ==
    \E c \in consumers : CheckReadable(c) \/ Read(c)
\end{lstlisting}

\section{Finals Tutorial Stuff}

\subsection{Go Concurrent Counter}
\begin{lstlisting}[language=Golang]
// worker
count := <-ch
count++
wg.Done()
ch <- count

// master
wg.Wait()
fmt.Println(<-ch)
\end{lstlisting}

\subsection{Go Multi Producer Multi Consumer Queue}
\begin{lstlisting}[language=Golang]
func producer(done chan struct{}, q chan int) {
    for {
        select {
        case q <- 1:
        case <-done:
            return
        }
    }
}

func consumer(done chan struct{}, q chan int, sumCh chan<- int) {
    sum := 0
    for {
        select {
        case <-done:
            sumCh <- sum
            close(sumCh)
            return
        case num := <-q:
            sum += num
        }
    }
}


var (
    NumProducer = 5
    NumConsumer = 5
)

func main() {
    start, done := make(chan struct{}), make(chan struct{})
    q := make(chan int)

    // 1 channer per consumer
    sumChs := make([]chan int, 0, NumConsumer)
    for i := 0; i < NumConsumer; i++ {
        sumCh := make(chan int, 1)
        sumChs = append(sumChs, sumCh)
    }

    for i := 0; i < NumProducer; i++ {
        go func() {
            <-start
            producer(done, q)
        }()
    }
    for j := 0; j < NumConsumer; j++ {
        j := j // capture j in the scope, for older Go versions
        go func() {
            <-start
            consumer(done, q, sumChs[j])
        }()
    }

    close(start) // start goroutines
    time.Sleep(time.Second) // run for 1 second
    close(done) // stop goroutines

    // collect all sums
    sum := 0
    for _, ch := range sumChs { // NOTE: range over slice, not combined channel
        sum += <-ch
    }
    fmt.Println("Sum: ", sum)
}
\end{lstlisting}

\subsection{Go Contexts}
\begin{lstlisting}[language=Golang]
// using a context
func producer(ctx context.Context, q queue) {
    for {
        select {
        case q<- 1:
        case <-ctx.Done(): return
        }
    }
}
// creating contexts
ctx, _ := context.WithTimeout(context.Background(), 4 * time.Second)
// context tree, if parent is canelled, child is also cancelled
ctx2, _ := context.withTimeout(ctx, 2 * time.Second)
// manually cancel a context
ctx3, cancel := context.WithCancel(ctx)
go func() { cancel() }()
\end{lstlisting}

\subsection{Go Fan-In Fan-Out}
\begin{lstlisting}[language=Golang]
type worker struct {
	inputCh  chan Event
	outputCh chan Event
}

func newWorker() *worker {
	return &worker{
		inputCh:  make(chan Event),
		outputCh: make(chan Event),
	}
}

// workers with promises
func (w *worker) start(
	ctx context.Context,
	fn EventFunc, workerQueue chan *worker, wg *sync.WaitGroup,
) {
    wg.Add(1)
	go func() {
        defer wg.Done()
        for {
            select {
            case <-ctx.Done():
                return
            case workerQueue <- w:
                e := <-w.inputCh
                w.outputCh <- fn(e)
            }
        }
	}()
}

// ordered multiplexer
func orderedMux(
	ctx context.Context, cancel context.CancelFunc,
	inputCh chan Event, workerQueue chan *worker, workerOutputCh chan chan Event,
) {
	go func() {
        for {
            select {
            case <-ctx.Done():
                return
            case e := <-inputCh:
                w := <-workerQueue
                w.inputCh <- e
                workerOutputCh <- w.outputCh
            }
        }
	}()
}

// in main
readerDone := make(chan struct{})
go func() {
    for outputCh := range workerOutputCh { // <-- reading a stream of promises
        output, more := <-outputCh // <-- read from the promise once only
        if !more {
            break // <-- a worker breaks its promise; workers exiting
        }
        fmt.Printf("Event id: %d\n", output.id)
    }
    close(readerDone)
}()

// serialiser
finalOutputCh := make(chan Event, 1)
	go func() {
        buffer := make(map[int64]Event)
        var nextEvent int64 = 1
        for output := range processedCh {
            buffer[output.id] = output
            for {
                if event, exists := buffer[nextEvent]; exists {
                    finalOutputCh <- event
                    nextEvent += 1
                } else {
                    break
                }
            }
        }
        close(finalOutputCh)
    }()
\end{lstlisting}

\subsection{Go Pipelining}
\begin{lstlisting}[language=Golang]
func toPipelineStage(
    ctx context.Context,
    reqCh <-chan Request, instances int,
    fn func(*Request) *Request,
) chan Request {
    outCh := make(chan Request)
    for i := 0; i < instances; i++ {
        go func() {
            for req := range reqCh {
                select {
                case outCh <- *fn(&req):
                case <-ctx.Done():
                    return
                }
            }
        }()
    }
    return outCh
}

// in main
c := make(chan Request, 100)
c = toPipelineStage(ctx, c, 2, step1)
c = toPipelineStage(ctx, c, 5, step2)
c = toPipelineStage(ctx, c, 10, step3)
c = toPipelineStage(ctx, c, 5, step4)
\end{lstlisting}

\subsection{C++ H2O using Barrier}
\begin{lstlisting}[language=C++]
std::counting_semaphore<> oxygenSem{1};
std::counting_semaphore<> hydrogenSem{2};
std::barrier<> barrier{3};

void oxygen(void (*bond)()) {
    oxygenSem.acquire();
    barrier.arrive_and_wait();
    bond();
    oxygenSem.release();
}

void hydrogen(void (*bond)()) {
    hydrogenSem.acquire();
    barrier.arrive_and_wait();
    bond();
    hydrogenSem.release();
}
\end{lstlisting}

\subsection{Go H2O using Daemon}
\begin{lstlisting}[language=Golang]
type WaterFactoryWithDaemon struct {
    // Channels for atoms to send their arrival requests
    precomH chan chan struct{}
    precomO chan chan struct{}
}

func NewFactoryWithDaemon() WaterFactoryWithDaemon {
    wfd := WaterFactoryWithDaemon{
        precomH: make(chan chan struct{}),
        precomO: make(chan chan struct{}),
    }

    // Daemon goroutine
    go func() {
        for {
            // Step 1: (Precommit)
            // Receive arrival requests from 2 hydrogen and 1 oxygen atoms
            h1 := <-wfd.precomH
            h2 := <-wfd.precomH
            o := <-wfd.precomO
            // Step 2: (Commit)
            // Tell the 3 atoms to start bonding
            h1 <- struct{}{}
            h2 <- struct{}{}
            o <- struct{}{}
            // Step 3: (Postcommit)
            // Wait until the 3 atoms have finished before looping
            // We re-use the same communication channel as (Commit)
            <-h1
            <-h2
            <-o
        }
    }()

    return wfd
}

func (wfd *WaterFactoryWithDaemon) hydrogen(bond func()) {
    commit := make(chan struct{}) // Step 1: Create private communication channel
    wfd.precomH <- commit         // Step 2: (Precommit)
    <-commit                      // Step 3: (Commit)
    bond()                        // Step 4: Bond
    commit <- struct{}{}          // Step 5: (Postcommit)
}

func (wfd *WaterFactoryWithDaemon) oxygen(bond func()) {
    commit := make(chan struct{}) // Step 1: Create private communication channel
    wfd.precomO <- commit         // Step 2: (Precommit)
    <-commit                      // Step 3: (Commit)
    bond()                        // Step 4: Bond
    commit <- struct{}{}          // Step 5: (Postcommit)
}
\end{lstlisting}

\subsection{Go H2O using Oxygen Leaders}
\begin{lstlisting}[language=Golang]
type WaterFactoryWithLeader struct {
    oxygenMutex chan struct{}
    precomH     chan chan struct{}
}

func NewFactoryWithLeader() WaterFactoryWithLeader {
    wf := WaterFactoryWithLeader{
        oxygenMutex: make(chan struct{}, 1),
        precomH:     make(chan chan struct{}),
    }
    wf.oxygenMutex <- struct{}{}
    return wf
}

func (wf *WaterFactoryWithLeader) hydrogen(bond func()) {
    commit := make(chan struct{}) // Step 1: Create private communication channel
    wf.precomH <- commit          // Step 2: (Precommit)
    <-commit                      // Step 3: (Commit)
    bond()                        // Step 4: Bond
    commit <- struct{}{}          // Step 5: (Postcommit)
}

func (wf *WaterFactoryWithLeader) oxygen(bond func()) {
    // Step 1: Become leader
    <-wf.oxygenMutex // For fun, we can use a channel as a mutex
    // Step 2: (Precommit)
    // Receive arrival requests from 2 hydrogen atoms
    h1 := <-wf.precomH
    h2 := <-wf.precomH
    // Step 3: (Commit)
    // Tell the 2 hydrogen atoms to start bonding
    h1 <- struct{}{}
    h2 <- struct{}{}
    // Step 4: Bond
    bond()
    // Step 5: (Postcommit)
    // Wait until the 2 hydrogen atoms to finish
    // We re-use the same communication channel as (Commit)
    <-h1
    <-h2
    // Step 6: Step down from being leader
    wf.oxygenMutex <- struct{}{}
}

// H2O2 problem
// other oxygen is follower who does not get oxygen mutex
case h2o2f.precomO <- commit: // Follower
    <-commit             // Step 3: (Commit)
    bond()               // Step 4: Bond
    commit <- struct{}{} // Step 5: (Postcommit)
\end{lstlisting}

\subsection{C++ FIFO Semaphore using Ticket Queue}
\begin{lstlisting}[language=C++]
struct FIFOSemaphore {
  std::mutex mut;
  std::condition_variable cond;
  std::atomic<std::ptrdiff_t> next_ticket;
  std::ptrdiff_t now_serving;

  FIFOSemaphore(std::ptrdiff_t initial_count)
      : mut{}, cond{}, next_ticket{1}, now_serving{initial_count} {}

  void acquire() {
    std::ptrdiff_t my_ticket =
        next_ticket.fetch_add(1, std::memory_order_relaxed);
    std::unique_lock lock{mut};
    cond.wait(lock, [=, this]() { return now_serving >= my_ticket; });
  }

  void release() {
    {
      std::scoped_lock lock{mut};
      now_serving++;
    }
    cond.notify_all();
  }
};
\end{lstlisting}

\subsection{C++ FIFO Semaphore using Queue of Semaphores}
\begin{lstlisting}[language=C++]
struct FIFOSemaphore {
  struct Waiter {
    std::binary_semaphore sem{0};
  };

  std::mutex mut;
  std::queue<std::shared_ptr<Waiter>> waiters;
  std::ptrdiff_t count;

  FIFOSemaphore(std::ptrdiff_t initial_count)
      : mut{}, waiters{}, count{initial_count} {}

  void acquire() {
    auto waiter = std::make_shared<Waiter>();
    {
      std::scoped_lock lock{mut};
      if (count > 0) {
        count--;  // Positive count,
        return;   // simply decrement without blocking
      }
      waiters.push(waiter);  // Zero count, add to waiters
    }
    waiter->sem.acquire();  // and block on the semaphore
  }

  void release() {
    std::shared_ptr<Waiter> waiter;
    {
      std::scoped_lock lock{mut};
      if (waiters.empty()) {
        count++;  // No waiters, simply increment count
        return;
      }

      waiter = waiters.front();  // Pop a waiter
      waiters.pop();
    }
    waiter->sem.release();  // and signal it
  }
};
\end{lstlisting}

\subsection{Go FIFO Semaphore using Buffered Channel}
\begin{lstlisting}[language=Golang]
type Semaphore struct {
	sem chan struct{}
}

func NewSemaphore(capacity int, initial_count int) *Semaphore {
	sem := Semaphore{
		sem: make(chan struct{}, capacity),
	}
	for i := 0; i < initial_count; i++ {
		sem.Release()
	}
	return &sem
}

func (s *Semaphore) Acquire() {
	// Receive from the channel to decrement the number of items
	// Blocks if there are no items (count is 0)
	<-s.sem
	// Blocked goroutines will be unblocked in FIFO order as of Go 1.24
}

func (s *Semaphore) Release() {
	// Send to the channel to increment the number of items
	s.sem <- struct{}{}
}
\end{lstlisting}

\subsection{Go FIFO Semaphore using Daemon}
\begin{lstlisting}[language=Golang]
type Semaphore struct {
	acquireCh chan chan struct{}
	releaseCh chan struct{}
}

func NewSemaphore(initial_count int) *Semaphore {
	sem := new(Semaphore)
	sem.acquireCh = make(chan chan struct{}, 100)
	sem.releaseCh = make(chan struct{}, 100)

	go func() {
		count := initial_count
		// The FIFO queue that stores the channels used to unblock waiters
		waiters := NewChanQueue()

		for {
			select {
			case <-sem.releaseCh: // Increment or unblock a waiter
				if waiters.Len() > 0 {
					ch := waiters.Pop()
					ch <- struct{}{} // Unblocks the oldest waiter
				} else {
					count++
				}

			case ch := <-sem.acquireCh: // Decrement or add a waiter
				if count > 0 {
					count--
					ch <- struct{}{} // Don't keep waiter blocked
				} else {
					waiters.PushBack(ch) // Add waiter to back of queue
				}
			}
		}
	}()

	return sem
}

func (s *Semaphore) Acquire() {
	ch := make(chan struct{})
	// Send daemon a channel that can be used to unblock us
	s.acquireCh <- ch
	// Block until daemon decides to unblock us
	<-ch
}

func (s *Semaphore) Release() {
	s.releaseCh <- struct{}{}
}
\end{lstlisting}

\subsection{Rust Concurrent Counter}
\begin{lstlisting}[language=Rust]
// mutex solution
use std::thread;
use std::sync::Mutex;

fn main() {
  let counter = Mutex::new(0);

  thread::scope(|s| {
    s.spawn(|| { 
        // THE FOLLOWING IS NOT IDIOMATIC RUST 
        // (what *counter.lock().unwrap() += 1 actually does)
        use std::ops::DerefMut;

        let lock_result = counter.lock();
        let mut lock_guard = lock_result.unwrap();
        let counter_ref: &mut i32 = lock_guard.deref_mut();
        *counter_ref += 1;
        // END UNIDIOMATIC RUST
    });
    s.spawn(|| { *counter.lock().unwrap() += 1; });
  });

  println!("{}", *counter.lock().unwrap());
}

// using Arc
use std::thread;
use std::sync::{Mutex, Arc};

fn main() {
  let counter = Arc::new(Mutex::new(0));

  let t0 = {
    let counter = counter.clone();
    thread::spawn(move || {
      // THE FOLLOWING IS NOT IDIOMATIC RUST
      use std::ops::{Deref, DerefMut};

      let mutex: &Mutex<i32> = counter.deref();
      let lock_result = mutex.lock();
      let mut lock_guard = lock_result.unwrap();
      let counter_ref: &mut i32 = lock_guard.deref_mut();
      *counter_ref += 1;
      // END UNIDIOMATIC RUST
    })
  };
  let t1 = {
    let counter = counter.clone();
    thread::spawn(move || {
      *counter.lock().unwrap() += 1;
    })
  };

  t0.join().unwrap();
  t1.join().unwrap();

  println!("{}", *counter.lock().unwrap());
}

// atomics
use std::sync::atomic::{AtomicI32, Ordering};
use std::thread;

use std::sync::{Mutex, Arc};

fn main() {
    let counter = AtomicI32::new(0);

    thread::scope(|s| {
        s.spawn(|| {
            counter.fetch_add(1, Ordering::Relaxed);
        });
        s.spawn(|| {
            counter.fetch_add(1, Ordering::Relaxed);
        });
    });

    println!("{}", counter.load(Ordering::Relaxed));
}
\end{lstlisting}

\subsection{Rust Line Echo Server}
\begin{lstlisting}[language=Rust]
use std::net::SocketAddr;

use tokio::io::{AsyncBufReadExt, AsyncWriteExt, BufReader};
use tokio::net::{TcpListener, TcpStream};

async fn handle_client(stream: TcpStream) -> std::io::Result<()> {
    let mut reader = BufReader::new(stream);
    let mut buf: Vec<u8> = Vec::new();
    loop {
        let size = reader.read_until(b'\n', &mut buf).await?;
        if size == 0 || buf[size - 1] != b'\n' {
            break;
        }
        reader.get_mut().write_all(&buf[..size]).await?;
        buf.clear();
    }
    Ok(())
}

#[tokio::main]
async fn main() -> std::io::Result<()> {
    let port = std::env::args()
        .nth(1)
        .map(|s| s.parse().unwrap())
        .unwrap_or(50000u16);
    let listener = TcpListener::bind(SocketAddr::from(([127, 0, 0, 1], port))).await?;
    loop {
        let (socket, _) = listener.accept().await?;
        tokio::spawn(async move {
            eprintln!("Accepted connection");
            std::mem::drop(handle_client(socket).await);
            eprintln!("Connection ended");
        });
    }
}
\end{lstlisting}

\subsection{Rust H2O}
\begin{lstlisting}[language=Rust]
#![allow(dead_code)]

use futures::future::join_all;
use std::sync::Arc;
use tokio::sync::{Barrier, Semaphore};

fn bond(s: &str) {
    println!("bond {s}!");
}

#[tokio::main]
async fn main() {
    let n = 100;
    let f = Arc::new(WaterFactory::new());

    let hs = (0..n * 2).map(|i| {
        let f = f.clone();
        tokio::spawn(async move {
            f.hydrogen(|| bond(&format!("h{i}"))).await;
        })
    });
    let os = (0..n).map(|i| {
        let f = f.clone();
        tokio::spawn(async move {
            f.oxygen(|| bond(&format!("o{i}"))).await;
        })
    });

    join_all(Iterator::chain(hs, os)).await;
}

struct WaterFactory {
    o_sem: Semaphore,
    h_sem: Semaphore,
    barrier: Barrier,
}

impl WaterFactory {
    fn new() -> Self {
        Self {
            o_sem: Semaphore::new(1),
            h_sem: Semaphore::new(2),
            barrier: Barrier::new(3),
        }
    }

    async fn oxygen(&self, bond: impl FnOnce()) {
        let _permit = self.o_sem.acquire().await.unwrap();
        self.barrier.wait().await;
        bond();
    }

    async fn hydrogen(&self, bond: impl FnOnce()) {
        let _permit = self.h_sem.acquire().await.unwrap();
        self.barrier.wait().await;
        bond();
    }
}
\end{lstlisting}

\section{PYP Questions}

\subsection{22/23 \texttt{io\_uring} in C++}
\begin{lstlisting}[language=C++]
struct Request {
    int client_fd;
    data_t data;
    res_t result;
    };

class ConcurrentRing {
private:
    LockFreeQueue<Request> queue;
    std::counting_semaphore write{SIZE};
    std::counting_semaphore read{0};
public:
    void submit_request(Request req) {
        write.acquire();
        queue.push(req);
        read.release();
    }
    Request retrieve_request() {
        read.acquire();
        std::optional<Request> req;
        while (true) {
            req = queue.try_pop();
            if (req) break;
        }
        write.release();
        return req.value();
    }
};

int main() {
    // Set up TCP socket
    int server_fd = socket(AF_INET, SOCK_STREAM, 0);
    // Bind socket to port
    listen(server_fd, 5);
    ConcurrentRing SQ, CQ;
    for (int i = 0; i < W; i++) {
        std::thread([&]() {
            while (true) {
                Request req = SQ.retrieve_req();
                req.process();
                CQ.submit_request(req);
            }
        });
        std::thread([&]() {
            while (true) {
                Request req = CQ.retrieve_req();
                send(req.client_fd, req);
            }
        });
    }
    while (true) {
        // Accept new client connection
        int client_fd = accept(server_fd, nullptr, nullptr);
        std::thread([&]() {
            data_t data = read(client_fd);
            while (data) {
                Request req{client_fd, data};
                SQ.submit_request(req);
                data = read(client_fd);
            }
        });
    }
}
\end{lstlisting}

\subsection{22/23 \texttt{io\_uring} in Go}
\begin{lstlisting}[language=Golang]
func main() {
    SQ = make(chan Request, SIZE)
    CQ = make(chan Request, SIZE)
    for i := 0; i < W; i++ {
        go func() {
            for {
                select {
                case req := <-SQ:
                    req.process()
                CQ <- req
                    default:
                }
            }
        }()
    }
    for i := 0; i < W; i++ {
        go func() {
            for {
                select {
                case req := <-CQ:
                    req.queue <- req
                default:
                }
            }
        }()
    }
    listener, err := net.Listen("tcp", ":8080")
    if err != nil {
        // handle error
    }
    for {
        conn, err := listener.Accept()
        if err != nil {
            // handle error
        }
        queue := make(chan Request)
        go func() {
            for {
                req := conn.read()
                req.queue = queue
                SQ <- req
            }
        }()
        go func() {
            for {
                select {
                case req := <-queue:
                    conn.send(req);
                default:
                }
            }
        }()
    }
}
\end{lstlisting}

\subsection{22/23 \texttt{io\_uring} in Rust}
\begin{lstlisting}[language=Rust]
async fn main() -> Result<(), Box<dyn Error>> {
    let addr = SocketAddr::from(([127, 0, 0, 1], 8080));
    let listener = TcpListener::bind(&addr).await?;
    let (sq_tx, mut sq_rx) = mpsc::channel::<IoOperation>(100);
    let (cq_tx, mut cq_rx) = mpsc::channel::<IoOperation>(100);
    loop {
        let (stream, _) = listener.accept().await?;
        tokio::spawn(async move {
            loop {
                let mut buf = Vec::new();
                stream.read(&mut buf).await;
                sq_tx.send(IoOperation::Read(stream, buf)).await;
            }
        });
    }
    tokio::spawn(async move {
        while let Some(op) = sq_rx.recv().await {
            tokio::spawn(async move {
                let res = process(op.data).await;
                cq_tx.send(IoOperation::Write(op.stream, res)).await;
            });
        }
    });
    tokio::spawn(async move {
        while let Some(op) = cq_rx.recv().await {
            op.stream.write_all(op.data).await;
        }
    });
    Ok(())
}
\end{lstlisting}

\subsection{24/25 Ticket Master}
\begin{lstlisting}[language=C++]
struct Seat {
    mutex mtx; // For per-seat locking
    atomic<bool> sold = false;
    int reservedBy = -1;
    steady_clock::time_point expiryTime = steady_clock::time_point::min();
};
class TicketSystem {
    vector<Seat> seats;
    thread expirationThread;
    atomic<bool> running = true;

    public:
    TicketSystem(int numSeats) : seats(numSeats) {
        expirationThread = thread(&TicketSystem::monitorExpirations, this);
    }
    ~TicketSystem() {
        running = false;
        if (expirationThread.joinable()) {
            expirationThread.join();
        }
    }
    bool processPayment(int userID, int seatID) {
        confirmPayment(userID);
        return true;
    }
    bool reserveSeat(int userID, int seatID, int holdSeconds = 5) {
        // (reservedBy + expiry for someone else holding seat)
        if (seatID < 0 || seatID >= seats.size()) return false;
        Seat& seat = seats[seatID];
        unique_lock<mutex> lock(seat.mtx);
        auto now = steady_clock::now();
        if (seat.sold || (seat.expiryTime > now && seat.reservedBy != userID)) {
            return false; // Already sold or held by someone
        }
        seat.reservedBy = userID;
        seat.expiryTime = now + seconds(holdSeconds);
        return true;
    }
    bool confirmSeat(int seatID, int userID) {
        if (seatID < 0 || seatID >= seats.size()) return false;
        Seat& seat = seats[seatID];
        unique_lock<mutex> lock(seat.mtx);
        auto now = steady_clock::now();
        if (seat.reservedBy == userID && seat.expiryTime > now) {
            seat.sold = true;
            seat.expiryTime = steady_clock::time_point::min();
            return true;
        } else {
            refundUser(userID);
            return false;
        }
    }
    void monitorExpirations() {
        while (true) {
            this_thread::sleep_for(seconds(1));
            auto now = steady_clock::now();
            for (auto& seat : seats) {
                unique_lock<mutex> lock(seat.mtx);
                if (!seat.sold && seat.expiryTime < now) {
                    seat.reservedBy = -1;
                    seat.expiryTime = steady_clock::time_point::min();
                }
            }
        }
    }
};

// atomics
struct Seat {
    atomic<bool> available;
    atomic<int> heldBy;
    Seat() : available(true), heldBy(0) {}
};
class TicketSystem {
    vector<Seat> seats;
    // ...
    bool reserveSeat(int userID, int seatID, int holdSeconds = 5) {
        steady_clock::time_point newExpiry =
        steady_clock::now() + seconds(holdSeconds);
        if (seatID < 0 || seatID >= seats.size()) return false;
        // Attempt to atomically reserve the seat
        bool expected = true;
        if (seats[seatID].available.compare_exchange_strong(expected, false)) {
            // Successfully changed availability to false
            seats[seatID].heldBy.store(userID);
            return true;
        } else {
            // Seat was not available
            return false;
        }
    }
};
\end{lstlisting}

\subsection{24/25 Load Balancer}
\begin{lstlisting}[language=Golang]
type Request struct {
    userID int
    seatID int
    responseCh chan bool
}
type Server struct {
    id int
    seatMin int
    seatMax int
    reqCh chan Request
    ts *TicketServer
}
type LoadBalancer struct {
    servers []*Server
}

func NewServer(id, seatMin, seatMax int) *Server {
    s := &Server{
        id: id,
        seatMin: seatMin,
        seatMax: seatMax,
        reqCh: make(chan Request),
        ts: NewTicketServer(seatMax),
        // assume Server handles seatMin to seatMax
    }
    go func() {
        for req := range s.reqCh {
            success := s.ts.HandleClient(req.userID,req.seatID)
            req.responseCh <- success
        }
    }()
    return s
}

func NewLoadBalancer(numServers int, maxSeatID int) *LoadBalancer {
    lb := &LoadBalancer{}
    seatsPerServer := (maxSeatID + numServers - 1) / numServers
    for i := 0; i < numServers; i++ {
        seatMin := i * seatsPerServer
        seatMax := seatMin + seatsPerServer - 1
        if seatMax > maxSeatID {
            seatMax = maxSeatID
        }
        lb.servers = append(lb.servers, NewServer(i, seatMin, seatMax))
    }
    return lb
}

func (lb *LoadBalancer) Dispatch(userID int, seatID int, responseCh chan bool) {
    for _, s := range lb.servers {
        if seatID >= s.seatMin && seatID <= s.seatMax {
            s.reqCh <- Request{
                userID: userID,
                seatID: seatID,
                responseCh: responseCh,
            }
            return
        }
    }
}

func main() {
    N := 3
    maxSeatID := 9
    lb := NewLoadBalancer(N, maxSeatID)
    results := make([]chan bool, 20)
    for i := 1; i <= 20; i++ {
        seatID := rand.Intn(10)
        results[i-1] = make(chan bool)
        lb.Dispatch(i, seatID, results[i-1])
    }
    for i := 1; i <= 20; i++ {
        confirmed := <-results[i-1]
        fmt.Printf("User %d seat confirmation: %v\n", i, confirmed)
    }
    time.Sleep(5 * time.Second)
}
\end{lstlisting}

\end{document}
